\documentclass{beamer}
\usetheme{Madrid}
\usepackage{pgfplots}
\pgfplotsset{compat=1.15}
\usepackage{mathrsfs}
\usetikzlibrary{arrows}
\usepackage[utf8]{inputenc}
\usepackage{xcolor}
\usepackage{nicematrix}
\usepackage{tikz}
\usepackage{comment}
\usepackage{hyperref}
\usepackage{fancyvrb}
\usepackage{mathtools}
\usepackage{comment}
\usepackage{array}
\usepackage{pgfcore}
\usepackage{listings}
\usepackage{color}
\definecolor{munsell}{rgb}{0.0, 0.5, 0.69}
\definecolor{dkgreen}{rgb}{0,0.6,0}
\definecolor{gray}{rgb}{0.5,0.5,0.5}
\definecolor{mauve}{rgb}{0.58,0,0.82}
\definecolor{minas}{RGB}{0.244, 0.172, 0.36}
\definecolor{PUJ}{RGB}{44, 86, 151}
\definecolor{PUJ2}{RGB}{43, 93, 156}
\definecolor{PUJ3}{RGB}{20, 52, 107}
\definecolor{cyandk}{rgb}{0.0, 0.72, 0.92}
\setbeamerfont{frametitle}{size=\LARGE ,series=\bfseries}
\setbeamercolor{frametitle}{fg=PUJ2, bg=white} %% title of the beamer
\setbeamercolor{titlelike}
{parent=structure,bg=PUJ2}
\setbeamercolor{title}{fg=white, bg=PUJ3} 
%\setbeamercolor{navigation symbols}{fg=white, bg=white}
\setbeamercolor*{palette primary}{use=structure,fg=black,bg=yellow}
\setbeamercolor*{palette secondary}{use=structure,fg=white,bg=PUJ3}
\setbeamercolor*{palette tertiary}{use=structure,fg=white,bg=PUJ3}

\setbeamercolor{block title}{bg=PUJ3,fg=white}


%% code information
\lstset{frame=tb,
  language=Python,
  aboveskip=3mm,
  belowskip=3mm,
  showstringspaces=false,
  columns=flexible,
  basicstyle={\small\ttfamily},
  numbers=none,
  classoffset=1,
  morekeywords={True,False}, keywordstyle=\color{munsell}, 
  classoffset=0, 
  keywordstyle=\color{blue},  
  commentstyle=\color{dkgreen},
  stringstyle=\color{PUJ3},
  numberstyle=\tiny\color{gray},
  breaklines=true,
  breakatwhitespace=true,
  tabsize=4,
}

%% You can change default language in the middle of document with \lstset{language=Java}.


%% PUT or Remove the logo in a slide.


%% Information topic

\institute{Javeriana}
\date{2020}

\title[Pontificia Universidad Javeriana] %optional
{Probability}
\subtitle{ using python.}

\author[
Iván Andrés Trujillo

] 
{
Iván Andrés Trujilllo Abella}

\institute[] 
{
  Facultad de Ingenieria\\
  Pontificia Universidad Javeriana
  \and
  
\textbf{}
}

\date[MINTA] % (optional)

\newif\ifplacelogo % create a new conditional
\placelogotrue % set it to true
%\logo{\ifplacelogo\color{red}\rule{.5cm}{.5cm}\fi}
\logo{\ifplacelogo \includegraphics[height= 2.0cm]{pujshield.eps}\fi}



\begin{document}



\frame{\titlepage}


\begin{frame}{Expected }
is defined as the mean of a random variable
\end{frame}


\begin{frame}{Variance}
\begin{equation}
\sigma^{2}  = E(x^{2}) - \mu^{2}
\end{equation}
remember that $\sum_{x}f(x)x = \mu$  and $\sum_{x} f(x) = 1 $
\begin{equation}
\begin{align*}
& \sigma^{2} = \sum_{x} (x - \mu)^{2}f(x)  \\
& = \sum(x^{2}f(x) - 2x\mu f(x) + \mu^{2}f(x))
\end{align*}
\end{equation}
Applying the algebra we find $\sum x^{2}f(x)-\mu^{2}$.

\end{frame}







\begin{frame}{Theorem 1}
$X$ is a random variable with a pdf $f(x)$ then $\mu$ of $g(x)$ is 
\begin{equation}
\mu_{g(x)} = E(g(x)) = \sum g(x)f(x)
\end{equation}
Example of income and the probability of sell a product.

\end{frame}



\begin{frame}{Theorem}
$X$ is a random variable with pdf $f(x)$ then then variance of $g(x)$ will be:

\begin{equation}
\sigma^{2}_{g(x)} = E( (g(x) - \mu_{g(x)})^{2}
\end{equation}
this is equation is derived of the definition of variance of a random variable, remember that $g(x)$ is a random variable with mean $\mu_{g(x)}$.
\end{frame}



\begin{frame}{Join distribution}

until now we try $\Omega$ in $\textbf{R}^{1}$ and we can be interested in find the probability of occurrence of two simultaneous random variables.
\begin{equation}
f(x,y) = P(X=x, Y=y)
\end{equation}
Some intuitive properties are:
\begin{itemize}
\item $f(x,y) \geq 0$
\item $\sum_{x}\sum_{y} f(x,y) =  1$ 
\end{itemize}

\begin{block}{Excercise}
Suppose the bag model with $n$ balls and there there are $r$ balls and $w$ balls where $r+w=n$ find the probability of get $x,y$ balls respectively.
\end{block}

\end{frame}

\begin{frame}{Marginal distribution}

\begin{equation}
g(x) = \sum_{y} f(x,y)
\end{equation}



\begin{equation}
h(y) = \sum_{x} = f(x,y)
\end{equation}
\end{frame}


\begin{frame}{Expected value of two random variables}
let be $X,Y$ two random variables with joint probability function distribution $f(x,y)$ the mean of $g(X,Y)$ is:
\begin{equation}
\mu_{g(X,Y)} = E[g(X,Y)] = \sum_{x} \sum_{y} g(x,y)f(x,y)
\end{equation}

\end{frame}



\begin{frame}{Covariance}{Insights}



\begin{equation}
\begin{align*}
\sigma_{X,Y} &= E[(X - \mu_{x})(Y - \mu_{y})] \\
 &= \sum_{x} \sum_{y} ( x  - \mu_{x}) (y - \mu_{y})f(x,y)
\end{align*}
\end{equation}

is a measure of association between two variables

\end{frame}



\begin{frame}{Shapiro wilk}


\end{frame}














\begin{frame}{P -values}

\end{frame}


\begin{comment}
%%%%%%%%%%%%%%%%%%%%%%%%%%%%%%%%%%%%%%%%%% WORDS
teardrop:
bonds:
heck:


\end{comment}

\begin{frame}{Hypothesis testing}
\textbf{Hypothesis} as idea or believe about a issue.

\end{frame}

\begin{frame}{Null hypothesis}
$H_{0}$ describe the current believe, and $H_{1}$ is a option if there is enough evidence to reject $H_{0}$.


\end{frame}


\begin{frame}{One tailed  test}
\begin{tikzpicture}[>={Stealth[length=6pt]},declare function={g(\x)=2*exp(-\x*\x/3);
    xmax=3.5;xmin=-3.4;x0=1.5;ymax=2.75;}]
 \draw[gray!50] (-3.7,0) edge[->] (4,0) foreach \X in {-3.5,-3,...,3}
  {(\X,0) -- ++ (0,0.1)} (0,0) edge[->] (0,ymax);
 \fill[gray!60] plot[domain=x0:xmax,samples=15,smooth] (\x,{g(\x)}) -- (xmax,0) -| cycle;  
 \draw[thick] plot[domain=xmin:xmax,samples=51,smooth] (\x,{g(\x)}); 
 \path (4,0) node[below]{$x$} (x0,0) node[below]{$x_{\alpha}$} 
  (0,ymax) node[right]{$f(x)$};
\end{tikzpicture}
\end{frame}





\begin{frame}{Two tailed test}{$H_{1} \neq value$}

\end{frame}



\begin{frame}{Test statistic}
test statistics is a value that allow us reject the null hypothesis, this uses the sampling statistics ( proportion, mean, or standard deviation) in a value of $z,t$ or $\chi^{2}$.3

\end{frame}



\begin{frame}{ANOVA test}


\end{frame}


\begin{frame}{t-test}

\end{frame}

\begin{frame}{$\chi^{2}$ test}

\end{frame}


\begin{frame}{Person}

\end{frame}


\begin{frame}{Spearman}

\end{frame}


\begin{frame}{Table one}
It is a useful algorithm to present working papers or give us insights about the problem.
\end{frame}





\end{document}